\section{Tangent Chern classes and the toric curves}
\label{sec:tangent-curves}

The boundary curves of the toric fiber give positive integration
currents on $Y$. Their Bott--Chern classes need not vanish with
their de Rham classes. We relate their sum to the second Chern
class of the tangent bundle, then compute each class by a period.
Positive exact combinations exclude a pluriclosed Hermitian metric.
The three curve classes have nonzero sum in Bott--Chern cohomology.

Retain the selected map $f:Y=Y^{\mathrm{ad}}(u)\to B=\mathbb P^1$,
its fixed sufficiently small nonzero $u$, and the divisor
$D=F_3+F_4+W$ of Section~\ref{sec:canonical-coherent}.
The integral increments, finite root fillings, adapted fan and
period units remain those of Theorem~\ref{cmp:thm:comparison}.
In particular, the following comparisons concern the ambient
cotangent bundle of this filled compact threefold.

\subsection{The three reduced curves}

Let $\nu:S\to W$ be the normalization, with
$S=\operatorname{Bl}_{3\text{ points}}\mathbb P^2$.
Write $B_1,\ldots,B_6$ for its cyclically ordered boundary curves,
and $C_j$ for the reduced image of $B_j$ and $B_{j+3}$,
for $1\le j\le3$.
Fix the indexing by the opposite ray pairs in the adapted angular
lattice:
\begin{equation}\label{tc:eq:curve-indexing}
 C_1:\ \pm(1,1),\qquad C_2:\ \pm(0,1),\qquad
 C_3:\ \pm(1,0).
\end{equation}
Let $p_+,p_-$ be the two triple points. We use the inclusions
\[
 i_j:C_j\hookrightarrow Y,\qquad
 k_\pm:\{p_\pm\}\hookrightarrow Y.
\]

\begin{lemma}\label{tc:lem:curve-embeddings}
Each restriction $B_j\to C_j$ and $B_{j+3}\to C_j$ is an
isomorphism, and $i_j$ is a smooth holomorphic embedding.
Thus $C_j\simeq\mathbb P^1$.
Each individual boundary curve has image degree one.
\end{lemma}

\begin{proof}
Use the height-one fan in Section~\ref{sec:toric-cusp}.
A boundary component in the star of the origin corresponds to
an edge $\{0,v\}$ of its triangulation.
A translation preserving this edge must either fix its endpoints
or exchange them. The first possibility is the zero translation.
The second would give $v=-v$, which is impossible for a nonzero
lattice vector.
Consequently no nonidentity period identifies two points in the
open part of this boundary curve.

The two triangles adjacent to an edge have different translation
types. Its two endpoints therefore map to the distinct triple
points $p_+$ and $p_-$.
Near either endpoint, the curve is a coordinate axis in the
smooth chart $q=x_0x_1x_2$.
The free period quotient is locally biholomorphic by
Proposition~\ref{c26:prop:quotient}.
It follows that the restriction to the entire compact boundary
curve is injective and locally an embedding. It is an isomorphism
onto its reduced image.
The opposite boundary component is a second normalization
component over that image, with the same degree one.

The simultaneous angular and deck shear
\eqref{hol:eq:cusp-shear}, followed by the specified unit
translation, identifies these charts with the selected adapted
filling. The period units occur in this biholomorphism and do not
change the degrees. The two translation classes of triangles give
exactly the two stated triple points.
\end{proof}

In particular, the two boundary currents together push forward
to $2[C_j]$, while either one pushes forward to $[C_j]$.
Here $[C_j]$ denotes integration over the reduced curve with its
complex orientation. There is no finite stabilizer factor.
At a triple point the third normalization branch is also present.
It is not an additional component dominating the same double curve.

\subsection{The coherent ambient cotangent identity}

Let $G_0(Y)$ be the Grothendieck group of coherent analytic
$\mathcal O_Y$-modules, with relations from short exact sequences.
Put
\[
 \mathcal L_0=f^*\mathscr L_B,\qquad
 \mathcal C_D=\bigoplus_{j=1}^3 i_{j*}\mathcal O_{C_j},\qquad
 \mathcal P_D=k_{+*}\mathcal O_{p_+}\oplus k_{-*}\mathcal O_{p_-}.
\]
Recall $\mathscr E$, $\mathcal K_{\mathrm{rel}}$ and
$\mathcal R_D=(\nu_D)_*\mathcal O_{S\sqcup F_3\sqcup F_4}$
from \eqref{nf:eq:relative-forms},
\eqref{nf:eq:scalar-sequence} and \eqref{hr:eq:residue-sequence}.

\begin{proposition}\label{tc:prop:coherent-cotangent}
The following identity holds in $G_0(Y)$:
\begin{equation}\label{tc:eq:coherent-cotangent}
 \begin{aligned}
 [\Omega_Y^1]
   ={}&[\mathcal L_0]+[\mathcal K_{\mathrm{rel}}]
                         +[\mathcal O_Y(-D)]\\
      &-\sum_{j=1}^3[i_{j*}\mathcal O_{C_j}]
                         +[k_{+*}\mathcal O_{p_+}]
                         +[k_{-*}\mathcal O_{p_-}].
 \end{aligned}
\end{equation}
The three curve terms have coefficient one.
\end{proposition}

\begin{proof}
The normalization complex \eqref{nf:eq:normalization-complex}
on $W$, together with the identity normalization on $F_3,F_4$,
gives the exact coherent complex on $Y$
\begin{equation}\label{tc:eq:conductor-complex}
 0\longrightarrow\mathcal O_D
 \longrightarrow\mathcal R_D
 \longrightarrow\mathcal C_D
 \longrightarrow\mathcal P_D\longrightarrow0.
\end{equation}
Its maps are the normalization unit and alternating branch
restrictions through the actual gluing maps.
Lemma~\ref{tc:lem:curve-embeddings} identifies the double-stratum
terms with the structure sheaves of the reduced embedded curves.

One can also see exactness directly on every stratum.
In a smooth local ring $A=\mathbb C\{x_1,x_2,x_3\}$, suppose
the reduced divisor has equation $x_1\cdots x_r=0$, with
$1\le r\le3$. A multiplying unit can be absorbed in one coordinate.
The local complex is
\begin{equation}\label{tc:eq:local-conductor}
 \begin{aligned}
 0\longrightarrow A/(x_1\cdots x_r)
  &\longrightarrow\bigoplus_a A/(x_a)\\
  &\longrightarrow\bigoplus_{a<b}A/(x_a,x_b)
   \longrightarrow\bigoplus_{a<b<c}A/(x_a,x_b,x_c)
   \longrightarrow0.
 \end{aligned}
\end{equation}
The intersection of the coordinate ideals is their product, so
the first map is injective. For a fixed surviving monomial,
its zero-exponent branches form a nonempty set.
The coefficient complex is the augmented simplex complex on
that set and is exact. Its contraction only copies or deletes
coefficients. It therefore preserves convergence on a common
smaller polydisk and proves exactness for analytic germs.

The two components dominating each $C_j$ are globally distinct
copies of $\mathbb P^1$. Ordering them defines their restriction
difference. At a triple point the other pair restrictions and
the point term retain the third branch.
The cyclic alternating relations give the last arrow.
Changing branch orders conjugates the complex by signs.
Thus the local complexes descend with the period units and give
\eqref{tc:eq:conductor-complex} on the filled $Y$.
Taking its class gives
\begin{equation}\label{tc:eq:normalization-class}
            [\mathcal R_D]=[\mathcal O_D]
                               +[\mathcal C_D]-[\mathcal P_D].
\end{equation}

The ambient residue sequence \eqref{hr:eq:residue-sequence}
has recipient $\mathcal R_D$.
For example, at a double point the residue of $dx_1/x_1$
is the independent branch tuple $(1,0)$.
With difference $f_2-f_1$, its restriction difference is $-1$.
It is not a matching tuple in $\mathcal O_D$.
This explains why replacing the residue recipient by
$\mathcal O_D$ would remove a necessary term.

The two bundle sequences
\eqref{nf:eq:absolute-log-sequence} and
\eqref{nf:eq:scalar-sequence} give
\[
 [\Omega_Y^1(\log D)]
       =[\mathcal L_0]+[\mathcal K_{\mathrm{rel}}]+[\mathcal O_Y].
\]
They use the actual base forms $3\,ds/s$ and $4\,ds/s$ at
the finite fibers, and the scalar logarithmic vector at the cusp.
Finally, the reduced divisor $D$ is Cartier, so
\[
 0\longrightarrow\mathcal O_Y(-D)
 \longrightarrow\mathcal O_Y\longrightarrow\mathcal O_D
 \longrightarrow0.
\]
Subtract \eqref{tc:eq:normalization-class} from the logarithmic
bundle identity and use this last sequence.
The result is \eqref{tc:eq:coherent-cotangent}.
\end{proof}

This identity concerns the locally free ambient sheaf $\Omega_Y^1$.
The intrinsic torsion-free cotangent sheaf on $W$ is a different
object, as in Theorem~\ref{ad:thm:ambient-kernel}.
The curve and point terms in \eqref{tc:eq:coherent-cotangent}
are present before taking characteristic classes.

\subsection{The normalized Bott--Chern character}

For a holomorphic vector bundle $E$ with a Hermitian metric,
let $R_E$ be the curvature of its Chern connection.
We use the normalization
\begin{equation}\label{tc:eq:character-normalization}
 c_{BC}(E)=\left[\det\left(1-\frac{R_E}{2\pi i}\right)\right]_{BC},
 \qquad
 \operatorname{ch}_{BC}(E)
     =\left[\operatorname{Tr}\exp\left(-\frac{R_E}{2\pi i}\right)
       \right]_{BC}.
\end{equation}
These classes are independent of the metric.
Their images in de Rham cohomology have the usual integral
Chern normalization.

We use two results about these classes on a compact complex
manifold \cite[Sections~2 and~8, Theorem~9.2]{BSWCoherent}.
First, the Chern character extends additively to $G_0$.
Second, for a smooth compact holomorphic embedding
$i:X\hookrightarrow Y$ with normal bundle $N_{X/Y}$,
\begin{equation}\label{tc:eq:embedding-RR}
 \operatorname{ch}_{BC}(i_*\mathcal O_X)
                 =i_*\operatorname{Td}_{BC}(N_{X/Y})^{-1}.
\end{equation}
The pushforward is duality of currents.
Its leading term is the integration current of $X$ with
coefficient one. The smooth and current definitions of
Bott--Chern cohomology agree by the elliptic decomposition
in \cite[Section~2.2]{BSWCoherent}.
These statements apply to coherent sheaves without requiring
a global finite resolution by holomorphic vector bundles.

Write $[C_j]_{BC}$ for the class of the integration current.
Let $\tau_{BC}\in H_{BC}^{3,3}(Y,\mathbb R)$ have integral one.
It is unique because integration is an isomorphism in this degree,
as proved in Proposition~\ref{bc:prop:products}.

\begin{theorem}\label{tc:thm:tangent-Chern}
For the selected threefold,
\begin{equation}\label{tc:eq:tangent-Chern}
 c_{2,BC}(T_Y)=\sum_{j=1}^3[C_j]_{BC},\qquad
 c_{3,BC}(T_Y)=2\tau_{BC}.
\end{equation}
The first class is represented by a nonzero positive current.
Its Bott--Chern class will be evaluated in
Theorem~\ref{tc:thm:evaluated-curves}.
\end{theorem}

\begin{proof}
The positive-degree ideal of the Bott--Chern algebra has square
zero by Theorem~\ref{bc:thm:ring}.
For every line bundle $L$ on $Y$, it follows that
\[
                     \operatorname{ch}_{BC}(L)=1+c_{1,BC}(L).
\]
The degree-$(2,2)$ part of \eqref{tc:eq:embedding-RR} is
\[
                \operatorname{ch}_{2,BC}(i_{j*}\mathcal O_{C_j})
                                   =[C_j]_{BC}.
\]
No additional factor of $2\pi i$ multiplies the current.
The normalization is already in the characteristic class.
Apply the additive character to \eqref{tc:eq:coherent-cotangent}.
Its degree-$(2,2)$ part gives
\[
 \operatorname{ch}_{2,BC}(\Omega_Y^1)
                         =-\sum_j[C_j]_{BC}.
\]
Since
$\operatorname{ch}_2(\Omega_Y^1)
=\tfrac12c_1(T_Y)^2-c_2(T_Y)$ and the square vanishes,
this proves the first identity in \eqref{tc:eq:tangent-Chern}.

For the top degree, $H^2(Y;\mathbb R)=0$ implies that
$\det T_Y$ has degree zero on every $C_j$.
The tangent-normal sequence on $C_j\simeq\mathbb P^1$ therefore
gives
\[
                     \deg N_{C_j/Y}=-2.
\]
No splitting type for this rank-two normal bundle is needed.
The linear term of its inverse Todd class is
$-\tfrac12c_1(N_{C_j/Y})$.
Equation~\eqref{tc:eq:embedding-RR} consequently gives
\[
 \int_Y\operatorname{ch}_{3,BC}(i_{j*}\mathcal O_{C_j})=1.
\]
Each such class equals $\tau_{BC}$, and each point class
$\operatorname{ch}_{3,BC}(k_{\pm*}\mathcal O_{p_\pm})$
also equals $\tau_{BC}$.
The top-degree part of \eqref{tc:eq:coherent-cotangent} is thus
\[
             \operatorname{ch}_{3,BC}(\Omega_Y^1)
                              =-3\tau_{BC}+2\tau_{BC}=-\tau_{BC}.
\]
All positive-degree products vanish, so the left side is
$-c_{3,BC}(T_Y)/2$. This proves the second identity.
Finally, $\sum_j[C_j]$ is a nonzero positive integration current
because the three curves are nonempty.
\end{proof}

The two point terms are essential in the top-degree comparison.
Keeping the three curve terms but omitting the point terms would
give $c_{3,BC}(T_Y)=6\tau_{BC}$.
Likewise, counting either boundary component with degree two
would double its reduced-cycle contribution before passing to
Bott--Chern cohomology.

\subsection{The invariant class on the actual complex tori}

To evaluate the curves, we first compare the integral fiber class
$\xi$ of Theorem~\ref{thm:geometric-level-six} with ordinary
Dolbeault cohomology on $Y$. The comparison has two parts: a
local holomorphic line bundle at the cusp, and an ordinary global
class with the same restriction. The finite logarithmic fillings
enter the second part.

Let $\delta$ be the positive line-fiber circle and let $\psi$ be
the primitive quotient covector of
Section~\ref{hol:sec:global-marking}.
The class $\xi$ satisfies
$\iota_\delta\xi=-6\psi$ and induces the positive unimodular
elliptic intersection form on $\ker\psi/\mathbb Z\delta$.

\begin{lemma}\label{tc:lem:invariant-type}
On every smooth fiber of the selected family, $\xi$ has type
$(1,1)$. It is the unique real $(1,1)$ alternating form with
this contraction and this restriction to $\ker\psi$.
\end{lemma}

\begin{proof}
Work on a simply connected disk in the regular base. Write the
elliptic curve as $\mathbb C/(\mathbb Z+\tau\mathbb Z)$, with
$\operatorname{Im}\tau>0$, and choose a holomorphic lift $\mu$
of the section $P$. The theta-ratio frame of
$\mathcal O(P-O)$ has multiplier $e^{2\pi i\mu}$ under the
$\tau$-period and multiplier one under the $1$-period.
The quotient generator covers translation by $-6P$.
Its multiplier is independent of the elliptic coordinate:
its ratio with a constant-multiplier lift is holomorphic on the
compact elliptic curve. Choose a logarithm $\sigma$ of this
multiplier, divided by $2\pi i$.
The actual additive periods are therefore
\begin{equation}\label{tc:eq:analytic-periods}
 \begin{gathered}
 \widetilde\alpha_1=(1,0),\qquad
 \widetilde\alpha_2=(0,1),\qquad
 \widetilde\beta_1=(\tau,\mu),\qquad
 \widetilde\beta_2=(-6\mu,\sigma),\\
 \Omega_{\mathrm{per}}=
       \begin{pmatrix}\tau&-6\mu\\ \mu&\sigma\end{pmatrix},
 \qquad \delta=\widetilde\alpha_2,
 \qquad \psi=\widetilde\beta_2^{\vee}.
 \end{gathered}
\end{equation}
These are the smooth-fiber counterparts of
\eqref{hol:eq:cusp-generators}. The entry $\sigma$ retains the
chosen linearization and the contraction parameter $u$.
No proportionality between $\sigma$ and $\tau$ is imposed.
The torus condition says that
$\operatorname{Im}\Omega_{\mathrm{per}}$ is invertible.

Put $Q_\xi=\operatorname{diag}(1,-6)$ and define an integral
real alternating form $E_\xi$ by
\begin{equation}\label{tc:eq:local-alternating-form}
 E_\xi(\widetilde\alpha_i,\widetilde\alpha_j)
 =E_\xi(\widetilde\beta_i,\widetilde\beta_j)=0,
 \qquad
 E_\xi(\widetilde\alpha_i,\widetilde\beta_j)=(Q_\xi)_{ij}.
\end{equation}
A real alternating form has type $(1,1)$ exactly when its
complexification vanishes on the kernel of the period map.
Vectors in that kernel are $(-\Omega_{\mathrm{per}}a,a)$.
Their pairing is
\[
 a^t\bigl(Q_\xi\Omega_{\mathrm{per}}
                  -\Omega_{\mathrm{per}}^tQ_\xi\bigr)a'.
\]
The matrix in parentheses is zero by
\eqref{tc:eq:analytic-periods}.
Thus $E_\xi$ has type $(1,1)$, the stated elliptic restriction,
and contraction $-6\psi$ with $\delta$.

The characterization is intrinsic. If $E$ and $E'$ have these
properties, their difference vanishes on
$\ker\psi\times\ker\psi$, so it is $\psi\wedge\ell$ for a
real covector $\ell$. Its contraction with $\delta$ is zero.
For a form of type $(1,1)$, this also makes its contraction
with $J\delta$ zero. Solving the imaginary part of the period
equations gives
\begin{equation}\label{tc:eq:complex-circle}
 \psi(J\delta)=
 \frac{\operatorname{Im}\tau}
      {\det(\operatorname{Im}\Omega_{\mathrm{per}})}\ne0.
\end{equation}
The identity
$(\psi\wedge\ell)(J\delta,-)=0$ now forces $\ell$ to be a
multiple of $\psi$. Hence $E=E'$.

The forms \eqref{tc:eq:local-alternating-form} consequently
agree under the actual changes of periods and frames.
Their coefficients are constant integers in each locally flat
period basis. They define a global flat integral class.
The rank-one invariant-lattice theorem and the contraction
normalization identify it with $\xi$, fixing its sign and scale.
Continuity of the Hodge subspaces extends this type statement
over each smooth good-reduction root family.
\end{proof}

The two signs in $Q_\xi$ are different. The type statement
therefore does not assert that $\xi$ is a polarization.
It uses the actual torsor multiplier and translation $-6P$.

\subsection{Local line bundles and an ordinary global lift}

Let $\Delta_\infty$ be a small disk about infinity and put
$N_\infty=f^{-1}(\Delta_\infty)$.
The toric retraction and integral specialization of
Theorem~\ref{c26:thm:integral} give
\begin{equation}\label{tc:eq:invariant-extension}
 H^2(N_\infty,\mathbb Z)\simeq H^2(W,\mathbb Z)
      \xrightarrow{\ p^*\ } H^2(F,\mathbb Z)^{T_\infty},
\end{equation}
where $F$ is a nearby smooth fiber and the second arrow is an
isomorphism. This is the comparison for the full period units,
transported by \eqref{hol:eq:cusp-isomorphism}.
Let $x_\infty$ be the unique integral extension of $\xi$.

\begin{proposition}\label{tc:prop:local-line}
There is a holomorphic line bundle $L_\infty$ on $N_\infty$
with $c_1(L_\infty)=x_\infty$.
There is also a nonzero section
\[
 \sigma_\xi\in H^0(B,\ker\kappa_1)
                 \subset H^0(B,\mathcal A_1)
\]
whose restriction to every smooth fiber is the class of $\xi$.
It is induced locally by closed real $(1,1)$-forms on the
absolute neighborhoods.
\end{proposition}

\begin{proof}
The sheaf $\mathcal C_2=R^2f_*\mathcal O_Y$ is locally free
at the cusp by Theorem~\ref{nf:thm:coherent} and its
base-change comparison. Leray over the Stein disk identifies
$H^2(N_\infty,\mathcal O)$ with its holomorphic sections.
The image of $x_\infty$ restricts to the $(0,2)$ component of
$\xi$ on every noncentral smooth fiber. That component is zero
by Lemma~\ref{tc:lem:invariant-type}.
A section of the locally free sheaf vanishing on the punctured
disk vanishes everywhere. The exponential sequence therefore
lifts $x_\infty$ to a holomorphic line bundle $L_\infty$.
Its normalized Chern curvature is a closed real $(1,1)$-form
representing $x_\infty$. In this normalization its integral on
a reduced compact curve is the integer degree of the line bundle.

At a finite exceptional point, pass to the smooth
good-reduction root disk. The pulled-back $\xi$ is a flat
integral class of type $(1,1)$ on every fiber.
Smooth proper base change and the same exponential-sequence
argument give a holomorphic line bundle $L$ on the root family.
For its actual free cyclic action of order $m=3$ or $4$, the
norm line
\[
                         \bigotimes_{g\in\mu_m}g^*L
\]
has the canonical permutation linearization and descends to
the quotient neighborhood $N_m$. Its pullback Chern class is
$m\xi$, since the flat class is invariant under this action.
Divide the normalized Chern curvature of the descended line by
$m$. This gives an absolute real $(1,1)$-class inducing $\xi$.
The division is in real or complex cohomology. It does not
assert integral descent of $\xi$ at order four; the integral
index of that filling is unchanged.

The same construction on regular disks gives local sections of
$\mathcal A_1=R^1f_*\mathscr E$.
They agree on regular overlaps, where they represent the same
fiber class $\xi$. Local freeness of $\mathcal A_1$ extends
that equality over the exceptional points.
They give the global section $\sigma_\xi$.
Each local section comes from absolute logarithmic one-form
cohomology, so its image under the connecting map $\kappa_1$
is zero. Its smooth-fiber restriction is nonzero.
\end{proof}

\begin{proposition}\label{tc:prop:ordinary-lift}
There is a class
\[
 z_{\bar\partial,\xi}=[\zeta_\xi]_{\bar\partial}
                      \in H^1(Y,\Omega_Y^1)
\]
mapping to $\sigma_\xi$. It is unique modulo the incoming
first Fr\"olicher differential. It satisfies
\begin{equation}\label{tc:eq:cusp-Dolbeault-comparison}
 z_{\bar\partial,\xi}|_{N_\infty}
       =c_1(L_\infty)\quad\text{in }H^1(N_\infty,\Omega^1),
 \qquad
 v_\xi=[\partial\zeta_\xi]_{\bar\partial}\ne0.
\end{equation}
Here $c_1(L_\infty)$ denotes the Dolbeault image of the
normalized Chern class, and $v_\xi$ is a cohomology class,
distinct from the contraction parameter $u$.
\end{proposition}

\begin{proof}
Use $\mathcal P_q$, $\mathcal K_1$ and
$\mathcal T=\mathbb C_0\oplus\mathbb C_1$ from
\eqref{hr:eq:absolute-direct-images}.
The exact sequence
\[
 0\longrightarrow\mathcal T\longrightarrow\mathcal P_1
        \longrightarrow\mathcal K_1\longrightarrow0,
 \qquad
 \mathcal K_1\simeq\mathcal O_B\oplus\mathcal O_B(-1)^{\oplus2}
\]
is onto on global sections. Since $H^1(B,\mathcal P_0)=0$,
Leray lifts $\sigma_\xi$ to
$H^1(Y,\Omega_Y^1(\log D))$.
The actual residue map identifies
\[
                 \mathcal T\xrightarrow{\ \sim\ }
                              R^1f_*\mathcal R_D.
\]
Its finite normalizations are
$dt/t=3\,ds/s$ and $d(t-1)/(t-1)=4\,ds/s$;
equation~\eqref{hr:eq:normalized-residue} gives the two inverse
generators. At the cusp its target is
$H^1(S,\mathcal O_S)=0$.
Subtract the unique torsion component having the same residue
as the logarithmic lift. The resulting class has residue zero.
The ordinary residue sequence \eqref{hr:eq:residue-sequence}
therefore supplies $z_{\bar\partial,\xi}$.

Its ambiguity is the degree-zero residue cokernel, a
one-dimensional line of divisor classes from the base.
These classes have $\partial$-closed Chern representatives.
The exact first-differential row in
Proposition~\ref{hr:prop:frolicher} identifies that line with
the incoming $d_1$ image.

It remains to compare the ordinary class at the cusp.
The map $\kappa_0$ is invertible on a sufficiently small
$\Delta_\infty$, including its full unit terms, by
Proposition~\ref{hr:prop:KS}.
The absolute logarithmic sequence on this Stein disk makes
\begin{equation}\label{tc:eq:first-cusp-injection}
 H^1(N_\infty,\Omega^1(\log W))
       \longrightarrow H^1(N_\infty,\mathscr E)
\end{equation}
injective: its kernel is the cokernel of $\kappa_0$ on
sections. Moreover,
$H^0(N_\infty,\mathcal R_W)=H^0(S,\mathcal O_S)=\mathbb C$,
and $d\log q$ has residue one. The residue sequence makes
\begin{equation}\label{tc:eq:second-cusp-injection}
 H^1(N_\infty,\Omega^1)
       \longrightarrow H^1(N_\infty,\Omega^1(\log W))
\end{equation}
injective as well.

The restrictions of $z_{\bar\partial,\xi}$ and of the
normalized Chern class of $L_\infty$ have the same image
$\sigma_\xi$ in the relative group.
The two injections prove the first equality in
\eqref{tc:eq:cusp-Dolbeault-comparison}.
The incoming-$d_1$ ambiguity has zero logarithmic image and
hence restricts to zero here.
Finally, the nonzero relative image of $z_{\bar\partial,\xi}$
shows that it is not on that incoming line.
Exactness of the same Fr\"olicher row gives $v_\xi\ne0$.
\end{proof}

\subsection{The three complex-oriented periods}

The local analytic basis in \eqref{tc:eq:analytic-periods}
has $\delta=\widetilde\alpha_2$.
The adapted cusp basis is different: the simultaneous
coordinate and deck shear \eqref{hol:eq:cusp-shear}, with the
chosen logarithm branches, gives
\begin{equation}\label{tc:eq:adapted-period-basis}
 \alpha_1=e_1-e_3,\qquad \alpha_2=e_3-e_2,\qquad
 \beta_1=e_3,\qquad \beta_2=e_4,
 \qquad \delta=-2\alpha_1+\alpha_2.
\end{equation}
All four periods enter the following computation.

\begin{theorem}\label{tc:thm:actual-periods}
With the indexing \eqref{tc:eq:curve-indexing} and the
normalized ordinary class of Proposition~\ref{tc:prop:ordinary-lift},
\begin{equation}\label{tc:eq:actual-periods}
 \int_{C_1}\zeta_\xi=-2,\qquad
 \int_{C_2}\zeta_\xi=3,\qquad
 \int_{C_3}\zeta_\xi=0.
\end{equation}
Each integral uses the reduced curve with its positive
complex orientation.
\end{theorem}

\begin{proof}
The matrix of $\xi$ in the common basis is
\[
 M_\xi=\begin{pmatrix}
 0&3&1&0\\-3&0&-2&0\\-1&2&0&-2\\0&0&2&0
 \end{pmatrix}.
\]
Conjugating by the columns of
\eqref{tc:eq:adapted-period-basis} gives the six coefficients
\begin{equation}\label{tc:eq:adapted-xi}
                 (0,1,2,2,-2,-2)
\end{equation}
in the lexicographic wedge order
$(\alpha_1\alpha_2,\alpha_1\beta_1,\alpha_1\beta_2,
  \alpha_2\beta_1,\alpha_2\beta_2,\beta_1\beta_2)$.
The degree-two specialization in \eqref{c26:eq:homologyp} is
\[
 P_2=\begin{pmatrix}
 0&-1&1&1&-1&0\\
 0&1&0&0&0&0\\
 0&0&0&0&1&0\\
 0&0&0&0&0&1
 \end{pmatrix}.
\]
Thus the central class $x_\infty$ has values
\begin{equation}\label{tc:eq:central-values}
       c=(2,3,0,-2)^t
       \quad\text{on }(E_0,E_1,E_2,h),
       \qquad P_2^tc=(0,1,2,2,-2,-2)^t.
\end{equation}
Uniqueness follows from the integral specialization theorem.

The cells $E_j$ have base-first orientations, which must now
be compared with the complex orientations.
The relevant vertices are
$v_0=(1,0)$, $v_1=(0,1)$, $v_2=(-1,1)$ and $v_3=(-1,0)$.
The oriented edges are
$e_0=[v_0,v_1]$, $e_1=[v_2,v_1]$, $e_2=[v_2,v_3]$,
with residual angle characters
$\lambda_0=(-1,1)$, $\lambda_1=(-1,0)$, $\lambda_2=(0,1)$.
Write $\zeta_{\mathrm{tor}}$ for the two dense torus coordinates.
For an edge from $a$ to $b$, the negative moment coordinate
\eqref{c26:eq:moment} has holomorphic residual coordinate
\[
 s=\zeta_{\mathrm{tor}}^{-(b-a)},\qquad
                  y=\frac{a+|s|^2b}{1+|s|^2}.
\]
Moving from $a$ to $b$ increases $|s|$.
The complex orientation is increasing $|s|$ followed by
increasing $\arg s$. The exponent $-(b-a)$ equals
$-\lambda_0$ on $e_0$, and equals $\lambda_1,\lambda_2$
on $e_1,e_2$. Consequently
\begin{equation}\label{tc:eq:complex-curve-cells}
                  [C_1]=-E_0,\qquad [C_2]=E_1,
                  \qquad [C_3]=E_2
\end{equation}
in $H_2(W,\mathbb Z)$.
The unit transport of Section~\ref{sec:toric-cusp} and the
adapted cusp isomorphism preserve each central stratum and
its complex orientation.

By \eqref{tc:eq:cusp-Dolbeault-comparison}, integration of
$\zeta_\xi$ on each curve equals the corresponding degree
of $L_\infty$.
Equations~\eqref{tc:eq:central-values} and
\eqref{tc:eq:complex-curve-cells} give $(-2,3,0)$.
Lemma~\ref{tc:lem:curve-embeddings} gives degree one for each
individual normalization branch, so no factor of two enters.
\end{proof}

\subsection{Period coordinates for the curve classes}

Let
\[
 j_{BC}:H_{BC}^{2,2}(Y)\longrightarrow H^2(Y,\Omega_Y^2)
\]
be the natural Dolbeault map.
The decomposition \eqref{bc:eq:22-lines} has
$L=\partial H_{\bar\partial}^{1,2}(Y)$.
The map $j_{BC}$ is injective on $L$ and zero on $\overline L$.
It is therefore injective on the real Bott--Chern plane:
if a real class has zero $L$ component, its conjugate component
also vanishes. This applies to current classes through the
smooth-current comparison above.

Use $\zeta=\zeta_\xi$ in Proposition~\ref{cq:prop:two-differentials},
and write
\[
       z=[\zeta]_A,\qquad v=v_\xi
                     =[\partial\zeta]_{\bar\partial}\ne0.
\]
Serre duality supplies a $\bar\partial$-closed form $\nu$ of
type $(1,2)$ with
\begin{equation}\label{tc:eq:period-normalization}
           \int_Y\partial\zeta\wedge\nu=1,
           \qquad w=[\partial\nu]_{BC}.
\end{equation}
The class $w$ generates $L$.
Define the three periods and their sum by
\begin{equation}\label{tc:eq:curve-periods}
                a_j=\int_{C_j}\zeta,\qquad A=a_1+a_2+a_3.
\end{equation}
The normalization of $v$ is now fixed by the integral class $\xi$.
The periods do not depend on its chosen $\bar\partial$-closed lift $\zeta$.

\begin{proposition}\label{tc:prop:curve-coordinates}
With these normalizations,
\begin{equation}\label{tc:eq:curve-coordinates}
 [C_j]_{BC}=-a_jw-\overline{a_j}\,\overline w,
 \qquad
 c_{2,BC}(T_Y)=-Aw-\overline A\,\overline w.
\end{equation}
For a real pluriclosed form $\omega$, define $t_C\in\mathbb C$ by
\[
           \iota\Phi([\omega]_A)=t_C v,
\]
using \eqref{cq:eq:forgetful} and \eqref{cq:eq:class-map}.
Then
\begin{equation}\label{tc:eq:Courant-area}
           \int_{C_j}\omega=\operatorname{Im}(t_C a_j),\qquad
           \int_Y\Bigl(\sum_j[C_j]\Bigr)\wedge\omega
                            =\operatorname{Im}(t_CA).
\end{equation}
\end{proposition}

\begin{proof}
Stokes' theorem and \eqref{tc:eq:period-normalization} give
\begin{equation}\label{tc:eq:period-pairings}
       \int_Y\zeta\wedge w=-1,
       \qquad \int_Y\zeta\wedge\overline w=0.
\end{equation}
The second equality uses $\bar\partial\zeta=0$ and
$\overline w=[\bar\partial\overline\nu]_{BC}$.
The two conjugate classes $w,\overline w$ form a basis.
Reality of $[C_j]_{BC}$ and the first pairing in
\eqref{tc:eq:period-pairings} determine its two coefficients,
giving the first formula in \eqref{tc:eq:curve-coordinates}.
The tangent formula follows by summing.

For independence of the lift, let $\zeta'$ be another such form
with $[\partial\zeta']_{\bar\partial}=v$.
The exact first-differential row of
Proposition~\ref{hr:prop:frolicher} gives
\[
         \zeta'-\zeta=\partial\gamma+\bar\partial\theta,
         \qquad \bar\partial\gamma=0,
\]
for forms $\gamma$ of type $(0,1)$ and $\theta$ of type $(1,0)$.
Their pullbacks to a compact curve have zero integral by Stokes,
since the other bidegree components vanish there.
Thus every $a_j$ is unchanged.

A real Aeppli class has the form $cz+\overline c\,\overline z$.
Theorem~\ref{cq:thm:class-map} gives $t_C=2ic$.
Its pairing with $[C_j]$ is
$ca_j+\overline c\,\overline{a_j}=2\operatorname{Re}(ca_j)$,
which equals $\operatorname{Im}(t_Ca_j)$.
This proves \eqref{tc:eq:Courant-area}.
\end{proof}

Let
\begin{equation}\label{tc:eq:curve-cone}
 \mathcal C_{\mathrm{tor}}
    =\left\{\sum_{j=1}^3\lambda_j[C_j]_{BC}:
                                  \lambda_j\in\mathbb R_{\ge0}\right\}.
\end{equation}
A nonzero nonnegative coefficient vector gives a relation with
zero precisely when
\begin{equation}\label{tc:eq:cone-zero}
              \lambda_1a_1+\lambda_2a_2+\lambda_3a_3=0.
\end{equation}
This tests Bott--Chern exactness of a linear combination of the
reduced curve currents. Its coefficients are real, independently
of integral divisor relations.

\begin{theorem}\label{tc:thm:evaluated-curves}
Put $\rho_{BC}=w+\overline w$. Then $\rho_{BC}\ne0$ and
\begin{equation}\label{tc:eq:evaluated-curves}
 [C_1]_{BC}=2\rho_{BC},\qquad
 [C_2]_{BC}=-3\rho_{BC},\qquad [C_3]_{BC}=0,
 \qquad c_{2,BC}(T_Y)=-\rho_{BC}\ne0.
\end{equation}
The complete nonnegative zero cone among these three reduced
curve classes is
\begin{equation}\label{tc:eq:exact-positive-cone}
 \left\{(\lambda_1,\lambda_2,\lambda_3)\in\mathbb R_{\ge0}^3:
                       \sum_j\lambda_j[C_j]_{BC}=0\right\}
       =\{(3s,2s,t):s,t\ge0\}.
\end{equation}
The sum of the three periods is $A=1$.
\end{theorem}

\begin{proof}
The two conjugate classes $w,\overline w$ are independent,
so their sum is nonzero.
Substitute \eqref{tc:eq:actual-periods} into
\eqref{tc:eq:curve-coordinates}.
The three curve classes are $(2,-3,0)\rho_{BC}$, and their
sum is $-\rho_{BC}$.
A real linear combination has zero class exactly when
$2\lambda_1-3\lambda_2=0$.
The nonnegative solutions are precisely
$(\lambda_1,\lambda_2,\lambda_3)=(3s,2s,t)$ with $s,t\ge0$.
Finally, $A=-2+3+0=1$.
\end{proof}

In particular, the positive relations do not make the sum
$[C_1]+[C_2]+[C_3]$ Bott--Chern exact.
For a real pluriclosed form the three curve integrals are
\begin{equation}\label{tc:eq:evaluated-areas}
 \left(\int_{C_1}\omega,\int_{C_2}\omega,\int_{C_3}\omega\right)
                   =(-2,3,0)\operatorname{Im}(t_C).
\end{equation}

\subsection{Global current potentials and the metric condition}

The current version of the real Bott--Chern quotient is
\begin{equation}\label{tc:eq:real-current-BC}
 H_{BC}^{2,2}(Y,\mathbb R)
  =\frac{\{T\text{ real of type }(2,2):dT=0\}}
         {\{i\partial\bar\partial S:
                               S\text{ real of type }(1,1)\}}.
\end{equation}
An auxiliary Hermitian metric gives a useful global primitive
whenever a class vanishes.
Put $P=\partial\bar\partial$, let $\Delta_{BC}$ be the real
elliptic Bott--Chern Laplacian, and let $G_{BC}$ be its Green
operator on the complement of its harmonic kernel.
We extend these operators to currents.

\begin{proposition}\label{tc:prop:global-potential}
For every real closed $(2,2)$-current $T$ with $[T]_{BC}=0$,
\begin{equation}\label{tc:eq:global-potential}
          S=-iP^*G_{BC}T
          \quad\text{is real and satisfies}\quad
          T=i\partial\bar\partial S.
\end{equation}
Consequently, a relation \eqref{tc:eq:cone-zero} with
$(\lambda_1,\lambda_2,\lambda_3)\ne0$ excludes a pluriclosed
Hermitian metric on $Y$.
In particular, $A=0$ excludes such a metric.
\end{proposition}

\begin{proof}
On the simultaneous kernel of $\partial$ and $\bar\partial$,
the Bott--Chern Laplacian is $PP^*$.
It preserves this closed subspace, as does its Green operator
\cite[equations~(2.9)--(2.13)]{BSWCoherent}.
The harmonic projection of $T$ vanishes exactly when
$[T]_{BC}=0$. Hence
\[
                        T=PP^*G_{BC}T.
\]
Complex conjugation anticommutes with $P$ and $P^*$ and
commutes with the real Green operator.
Thus $P^*G_{BC}T$ is purely imaginary, and
\eqref{tc:eq:global-potential} follows.
The smooth-current elliptic decomposition supplies this identity
and closed range on the entire compact $Y$.
No local potentials need to be patched across finite or cusp
overlaps.

For a nonzero nonnegative vector $\lambda$, the current
$T_\lambda=\sum_j\lambda_j[C_j]$ is positive and nonzero.
If \eqref{tc:eq:cone-zero} holds, it has the global potential
just constructed. A strictly positive pluriclosed form $\omega$
would satisfy
\[
        0<\sum_j\lambda_j\int_{C_j}\omega
          =\langle i\partial\bar\partial S,\omega\rangle
          =\langle S,i\partial\bar\partial\omega\rangle=0,
\]
a contradiction. For $A=0$, take all three coefficients equal
to one.
\end{proof}

There is also an unconditional transgression formula.
If $c_2(T_Y,h)$ is a smooth Chern--Weil representative in
\eqref{tc:eq:character-normalization}, then
\[
 U=\sum_j[C_j]-c_2(T_Y,h)
\]
has zero Bott--Chern class by Theorem~\ref{tc:thm:tangent-Chern}.
Applying \eqref{tc:eq:global-potential} to $U$ gives a real
global current whose $i\partial\bar\partial$ is this difference.
This formula does not require $A=0$.

Define the cone of classes of pluriclosed Hermitian metrics by
\begin{equation}\label{tc:eq:SKT-cone}
 \mathcal P_{\mathrm{SKT}}(Y)
 =\{[\omega]_A:\omega>0,\ \overline\omega=\omega,
                         \ \partial\bar\partial\omega=0\}
\end{equation}
where $\omega$ is smooth.

\begin{theorem}\label{tc:thm:no-SKT}
The selected threefold $Y$ admits no pluriclosed Hermitian
metric:
\[
                         \mathcal P_{\mathrm{SKT}}(Y)=\varnothing.
\]
For $s,t\ge0$, the global current
\[
                  T_{s,t}=s(3[C_1]+2[C_2])+t[C_3]
\]
has a real potential
$S_{s,t}=-iP^*G_{BC}T_{s,t}$ satisfying
$T_{s,t}=i\partial\bar\partial S_{s,t}$.
It is positive and nonzero whenever $(s,t)\ne(0,0)$.
\end{theorem}

\begin{proof}
Theorem~\ref{tc:thm:evaluated-curves} makes the harmonic
Bott--Chern projection of $T_{s,t}$ zero.
Proposition~\ref{tc:prop:global-potential} gives its stated
potential on the entire compact $Y$.
The coefficients are nonnegative and the curves are distinct
nonempty reduced curves, so their currents cannot cancel.
If $\omega$ were a pluriclosed Hermitian form, then
\[
 0<\int_{C_3}\omega
    =\langle i\partial\bar\partial S_{0,1},\omega\rangle
    =0
\]
by integration by parts. This is a contradiction.
The other generator $3[C_1]+2[C_2]$ gives a second positive
exact current.
\end{proof}

The auxiliary Hermitian metric used to define $G_{BC}$ is
unrestricted. The obstruction excludes every smooth positive
pluriclosed form, not only forms invariant under a chosen action.
The class map of Theorem~\ref{cq:thm:class-map} and its
holomorphic Courant algebroids remain defined on all real
Aeppli classes; none of those classes has a pluriclosed
Hermitian representative on this $Y$.
The fixed nonzero parameter and selected finite and cusp
fillings remain part of this assertion.

\subsection{A Calabi--Eckmann comparison}

The distinction between de Rham and Bott--Chern classes is
visible on a different compact threefold.
Consider
\begin{equation}\label{tc:eq:CE-quotient}
 M_{CE}=
 \bigl((\mathbb C^2\setminus\{0\})\times
       (\mathbb C^2\setminus\{0\})\bigr)/\mathbb C,
 \qquad
 \gamma\cdot(z,w)=(e^\gamma z,e^{i\gamma}w).
\end{equation}
This is the Calabi--Eckmann example of
\cite[Example~2.8]{GFJSCourant}.
Writing $\gamma=s+it$, the two norms change by $e^s$ and
$e^{-t}$. Every orbit therefore has a unique representative
with $|z|=|w|=1$. The quotient is diffeomorphic to
$S^3\times S^3$ and fibers holomorphically over
$\mathbb P^1\times\mathbb P^1$ with elliptic fibers.

On the two Hopf factors choose the real connections $\mu_j$
with circle period $2\pi$, and write $\nu_j=d\mu_j$.
The forms $\nu_j$ are pullbacks of positive base forms with
integral $2\pi$.
Let $\xi_j$ be the circle fields, normalized by
$\mu_j(\xi_j)=1$.
The quotient complex structure has $J\xi_1=\xi_2$ and hence
\[
 J^*\mu_1=-\mu_2,\qquad J^*\mu_2=\mu_1,
 \qquad J^*\nu_j=\nu_j.
\]
The horizontal complex structures are those of the two bases.
Thus
\begin{equation}\label{tc:eq:CE-metric}
                  \omega_{CE}=\nu_1+\nu_2+\mu_1\wedge\mu_2
\end{equation}
is a positive real $(1,1)$-form. In particular,
$\omega_{CE}(\xi_1,J\xi_1)=1$.
Its torsion, for $d^c=i(\bar\partial-\partial)$, is
\begin{equation}\label{tc:eq:CE-torsion}
 \begin{aligned}
 d\omega_{CE}&=\nu_1\wedge\mu_2-\mu_1\wedge\nu_2,\\
 H_{CE}=-d^c\omega_{CE}=J^*d\omega_{CE}
               &=\nu_1\wedge\mu_1+\nu_2\wedge\mu_2,\\
 dH_{CE}&=\nu_1^2+\nu_2^2=0.
 \end{aligned}
\end{equation}
Both torsion signs are positive for the stated dual-connection
convention. Each square vanishes because its form comes from
a curve. Therefore \eqref{tc:eq:CE-metric} is pluriclosed.

On an elliptic fiber $E$, the complex orientation is
$d\theta_1\wedge d\theta_2$, and
\begin{equation}\label{tc:eq:CE-area}
 \int_E\omega_{CE}
    =\int_0^{2\pi}\int_0^{2\pi}d\theta_1\,d\theta_2
    =4\pi^2.
\end{equation}
Yet $H_2(M_{CE};\mathbb R)=H^4(M_{CE};\mathbb R)=0$.
Thus $[E]$ is de Rham exact as a current, while
$[E]_{BC}\ne0$.
If $[E]=i\partial\bar\partial S$, pairing with the same
$\omega_{CE}$ would give $0=4\pi^2$ by integration by parts.
The positive area is therefore a direct obstruction to that
claimed Bott--Chern primitive.

The two torsion periods also have their fixed values:
\[
 \int_{S^3\times\{*\}}H_{CE}=4\pi^2,
 \qquad
 \int_{\{*\}\times S^3}H_{CE}=4\pi^2,
\]
with the Hopf orientations. Hence $[H_{CE}]$ is nonzero in
$H^3(M_{CE};\mathbb R)$.
On the selected $Y$, by contrast, $H^3(Y;\mathbb R)=0$ and
every closed real torsion form is exact.
This difference does not transfer the Calabi--Eckmann metric
to $Y$. On the selected $Y$, the positive exact curve current
of Theorem~\ref{tc:thm:no-SKT} excludes such a metric.
